\newcounter{joparagraphcounter}
\setcounter{joparagraphcounter}{1}

\section{Jugendordung}

\subsection{Präambel}
Die Jugend des Narrenbund Ostfildern e.V. ist Teil der närrischen Gemeinschaft und verfolgt das Ziel, junge Menschen an das Brauchtum und die Tradition der Fasnet heranzuführen. Diese Jugendordnung regelt die Aufgaben, Rechte und Strukturen der Vereinsjugend.

\subsection{§ \thejoparagraphcounter\ Erlass, Änderung, Aufhebung und Bekanntmachung} \label{jo:par:erlass_aenderung_aufhebung}
\begin{enumerate}
    \item \label{jo:abs:of1a22} Diese Jugendordnung kann jederzeit durch den Vereinsausschuss erarbeitet und beschlossen werden.

    \item \label{jo:abs:rqn6xb} Änderungen werden durch den Beschluss des Vereinsausschusses wirksam und allen Mitgliedern bekannt gemacht.
\end{enumerate}

\refstepcounter{joparagraphcounter}
\subsection{§ \thejoparagraphcounter\ Zugehörigkeit zur Jugend} \label{jo:par:zugehoerigkeit_jugend}
\begin{enumerate}
    \item \label{jo:abs:s1rqff} Zur Jugend des Vereins gehören alle Mitglieder bis einschließlich des 27. Lebensjahres.

    \item \label{jo:abs:eh777b} Die Jugend ist organisatorisch vollständig in den Verein integriert und bildet keine eigenständige Untergliederung.
\end{enumerate}

\refstepcounter{joparagraphcounter}
\subsection{§ \thejoparagraphcounter\ Jugendleitung} \label{jo:par:jugendleitung}
\begin{enumerate}
    \item \label{jo:abs:fnqb6q} Die Jugendleitung besteht aus:
    \begin{enumerate}
        \item Jugendleiter
        
        \item Stellv. Jugendleiter
    \end{enumerate}

    \item \label{jo:abs:86krfo} Die Jugendleitung wird von den Mitgliedern der Jugend gewählt; diese Wahl sollte zeitlich vor der Mitgliederversammlung des Vereins liegen.

    \item \label{jo:abs:0hhncw} Die Wahl bedarf in jedem Fall der Bestätigung durch die Mitgliederversammlung.

    \item \label{jo:abs:bq5ltj} Die Amtszeit beträgt zwei Jahre; Wiederwahl ist möglich.
\end{enumerate}

\refstepcounter{joparagraphcounter}
\subsection{§ \thejoparagraphcounter\ Rechte der Jugendleitung im Verein} \label{jo:par:rechte_jugendleitung}
\begin{enumerate}
    \item \label{jo:abs:oso8pe} Der Jugendleiter hat gemäß \parabsref{sa:par:vereinsausschuss}{sa:abs:g0t862} der Satzung Sitz und Stimme im Ausschuss des Vereins.

    \item \label{jo:abs:xo7r87} Der stellv. Jugendleiter nimmt nur im Verhinderungsfall des Jugendleiters an Ausschusssitzungen teil und hat dort kein Stimmrecht.

    \item \label{jo:abs:onf6im} Die Jugendleitung ist berechtigt, dem Ausschuss Vorschläge zur Jugendarbeit zu unterbreiten.
\end{enumerate}

\refstepcounter{joparagraphcounter}
\subsection{§ \thejoparagraphcounter\ Haftung und Verantwortung} \label{jo:par:haftung_verantwortung}
\begin{enumerate}
    \item \label{jo:abs:gh3w5l} Die Jugendleitung haftet gegenüber dem Verein und seinen Mitgliedern nur für vorsätzliche oder grob fahrlässige Schäden, im Rahmen der gesetzlichen Ehrenamtshaftung.

    \item \label{jo:abs:up1q4u} Die Jugendleitung trägt keine volle persönliche Haftung für Tätigkeiten oder Entscheidungen der Jugend, sondern unterstützt diese im Rahmen ihrer Möglichkeiten.

    \item \label{jo:abs:av3mbu} Der Verein unterhält die üblichen Vereinsversicherungen (Haftpflicht/Unfall/etc.), in deren Rahmen die Tätigkeit der Jugendleitung mit abgedeckt ist, soweit gesetzlich möglich.
\end{enumerate}

\refstepcounter{joparagraphcounter}
\subsection{§ \thejoparagraphcounter\ Inkrafttreten} \label{jo:par:inkrafttreten}
\begin{enumerate}
    \item \label{jo:abs:x95wtg} Diese Jugendordnung tritt nach Beschluss der Ausschusssitzung vom dd. MMMM YYYY in Kraft. \\
    Sie ersetzt alle vorherigen Versionen
\end{enumerate}